\ncdmodultest{1}{Fritid, Hverdag og Transport}{全真模考卷一：夜校课程、通勤故事与二手文化}{成人夜校与业余课程 / 城铁遗落背包 / 丹麦二手与跳蚤市场文化 / 三人谈通勤方式 / 生日聚会回信与夜校请假 / 工作与住房口语陈述及幼儿园信息差}{现在完成时 vs 一般过去时切换 · 从句否定与副词前置 · 篇章连接副词（derfor/alligevel/desuden）· 关系代词 som/der · 礼貌半正式请假表达}

% =====================================================================
% DEL 1: LÆSNING
% =====================================================================
\ncdsubsec{Del 1: Læsning（阅读理解｜40 minutter · 24 point）}

\begin{tipbox}{考场规则与配速建议（Læsning｜16--24 point = Bestået）}
\small
本部分共四道大题，总计 \textbf{24 分}，及格线为 \textbf{16 分}。考场限时 \textbf{40 分钟}，全程\textbf{禁用任何词典或手机}。建议配速：Opgave 1（8 分钟）$\rightarrow$ Opgave 2（8 分钟）$\rightarrow$ Opgave 3（8 分钟）$\rightarrow$ Opgave 4（12 分钟），最后留 4 分钟复查。
\end{tipbox}

% ---------------------------------------------------------------------
\subsubsection*{Opgave 1 (NY · 6 point)：Aftenskolekurser i byen}

\textbf{Instruktion:}
\begin{itemize}
  \item Læs de seks små tekster om personerne, der gerne vil gå til et kursus i deres fritid (1--6).
  \item Find den annonce (B--I), der passer til hver tekst.
  \item Skriv bogstavet på annoncen, der passer til teksten. Der er to annoncer, du ikke skal bruge. Se eksemplet (0).
\end{itemize}

\begin{enumerate}
  \item[(0)] Jonas vil gerne lave noget kreativt med sine hænder, men han arbejder sent hver aften fra mandag til fredag og har kun tid i weekenden. Han drømmer om at lære at forme kopper og skåle i ler, selvom han aldrig har prøvet det før.\\
  \textbf{Svar:} \underline{\textbf{\ A\ }} \quad \textcolor{NCDGray}{\small [Eksempel: Svar A]}

  \item[1.] Astrid på 62 år har lige købt en ny telefon og vil gerne blive bedre til at tage skarpe billeder af sine børnebørn og sine ferier. Hun har ikke et stort kamera og ønsker ikke et langt kursus, men leder efter et kort aftenhold for begyndere, hvor man lærer at bruge telefonens kamera og redigere fotos.\\
  \textbf{Svar:} \rule{1.2cm}{0.4pt}

  \item[2.] Tobias elsker at stå i køkkenet, men han vil gerne spise mindre kød i hverdagen. Han savner nye opskrifter på varme grøntsagsretter med gode krydderier. Det er vigtigt for ham, at holdet laver maden sammen i et hyggeligt fællesskab og spiser aftensmaden bagefter. Han kan kun om onsdagen eller torsdagen efter arbejde.\\
  \textbf{Svar:} \rule{1.2cm}{0.4pt}

  \item[3.] Yasmin skal rejse tre uger rundt i Spanien til sommer og vil gerne kunne tale lidt med de lokale. Hun har aldrig lært spansk før og er ikke interesseret i svær grammatik eller skriftlige prøver. Hun søger et begynderhold, hvor man øver sig i at snakke, bestille mad på café og spørge om vej.\\
  \textbf{Svar:} \rule{1.2cm}{0.4pt}

  \item[4.] Anders cykler tolv kilometer til arbejde hver dag og bruger alt for mange penge hos cykelsmeden, når dækket punkterer, eller bremserne ikke virker godt. Han leder efter et praktisk værksted om aftenen, hvor han kan tage sin egen cykel med og selv lære at lappe slanger, rense kæden og justere gear.\\
  \textbf{Svar:} \rule{1.2cm}{0.4pt}

  \item[5.] Birgitte elsker musik og savner at synge sammen med andre efter en lang arbejdsdag. Hun kan desværre ikke læse noder og vil ikke til en svær optagelsesprøve eller synge alene foran andre. Hun leder efter et afslappet kor om aftenen, hvor man hygger sig med kendte danske sange og popsange.\\
  \textbf{Svar:} \rule{1.2cm}{0.4pt}

  \item[6.] Omar arbejder foran en skærm otte timer om dagen og får ofte ondt i nakken og ryggen. Han ønsker ikke hård fitness med høj musik eller konkurrence, men vil gerne gå til et roligt hold om aftenen, hvor man laver blide strækøvelser, træner vejrtrækning og får mere ro i kroppen og hovedet.\\
  \textbf{Svar:} \rule{1.2cm}{0.4pt}
\end{enumerate}

\vspace{0.3em}
\small
\renewcommand{\arraystretch}{1.35}
\begin{longtable}{@{} >{\raggedright\arraybackslash}p{8.2cm} @{\hspace{0.6cm}} >{\raggedright\arraybackslash}p{8.2cm} @{}}
\toprule
\multicolumn{2}{@{} >{\raggedright\arraybackslash}p{17.0cm} @{}}{%
  \badge{NCDRed}{A}\ \textbf{\color{NCDNavy} Keramik i weekenden -- Aarhus Keramikværksted} \quad \textcolor{NCDGray}{\small [Eksempel]}\newline
  Har du lyst til at arbejde med hænderne i rolige rammer om lørdagen fra kl. 10.00 til 13.00? På vores weekendhold lærer du at forme smukke kopper, tallerkener og skåle i blødt ler og dekorere dem med flotte glasurer. Du behøver ingen tidligere erfaring, og alle materialer samt brænding i vores store ovn er med i prisen. Tilmeld dig på keramik-aarhus.dk.%
} \\
\midrule
\endfirsthead
\toprule
\multicolumn{2}{@{} >{\raggedright\arraybackslash}p{17.0cm} @{}}{%
  \small\color{NCDGray}\textit{Opgave 1 · Aftenskolekurser i byen（续表）}%
} \\
\midrule
\endhead
\bottomrule
\endfoot
\badge{NCDNavy}{B}\ \textbf{\color{NCDNavy} Alle kan synge! Popkor uden noder}\newline
Drømmer du om at synge i kor, men tror du, at man skal kunne læse noder? Hos Klang \& Glæde handler det først og fremmest om sangglæde og fællesskab! Hver tirsdag kl. 19.00--21.00 synger vi kendte danske klassikere og nyere pophits sammen i et stort kor. Ingen synger solo, og der er ingen optagelsesprøve. Efter pausen drikker vi kaffe sammen. Se mere på klangogglaede.dk.
&
\badge{NCDNavy}{F}\ \textbf{\color{NCDNavy} Surdejsbrød og franske kager}\newline
Lær hemmeligheden bag sprødt surdejsbrød, croissanter og klassiske lagkager! Vi mødes hver søndag formiddag kl. 10.00--14.00 i vores bageri på Havnegade. Underviseren viser trin for trin, hvordan man ælter dej og tempererer chokolade. Du får selvfølgelig alle dine nybagte brød og kager med hjem i en pose til familien. Holdet koster 950 kr. for fire søndage. Tilmelding: bagervaerk.dk.
\\
\midrule
\badge{NCDNavy}{C}\ \textbf{\color{NCDNavy} Grøn hverdagsmad -- Madskolen i Midtbyen}\newline
Har du lyst til at lave spændende retter uden kød, som smager af meget? Hver onsdag fra kl. 17.30 til 20.30 mødes vi i vores store skolekøkken og tilbereder nemme, varme retter med bønner, linser, friske grøntsager og aromatiske krydderier. Vi arbejder i små grupper og slutter altid aftenen med at spise ved et langt bord sammen. Råvarer er med i prisen. Læs mere på madskolen-midt.dk.
&
\badge{NCDNavy}{G}\ \textbf{\color{NCDNavy} Fotografering med mobilen -- Kulturhuset}\newline
Vil du lære at tage flotte og tydelige billeder uden et dyrt spejlreflekskamera? På dette korte kursus over fire tirsdag aftener fra kl. 18.30 til 20.30 lærer du enkle tricks til lys, portrætter og feriebilleder på din smartphone. Vi viser dig også gratis apps til hurtig billedredigering. Holdet er for alle uden erfaring, og du skal kun medbringe din egen mobiltelefon. Tilmelding på kulturhus-aften.dk.
\\
\midrule
\badge{NCDNavy}{D}\ \textbf{\color{NCDNavy} Naturfoto med systemkamera -- Østjysk Fotoklub}\newline
Har du allerede et godt spejlreflekskamera, og kender du de grundlæggende indstillinger? På dette 12-ugers kursus tager vi ud i skoven og ved kysten hver lørdag formiddag kl. 9.00--13.00 for at fotografere fugle og landskaber i al slags vejr. Du skal selv medbringe digitalt systemkamera med udskiftelige objektiver og stativ. Holdet kræver grundlæggende erfaring med manuel fotografering. Info: fotoklub-oest.dk.
&
\badge{NCDNavy}{H}\ \textbf{\color{NCDNavy} Rolig yoga og udstrækning -- Aftenro Bevægelse}\newline
Trænger din krop til at komme ned i tempo efter en travl dag på kontor eller studie? Mandag og onsdag aften kl. 19.30--20.45 tilbyder vi blid yoga i et stille, varmt lokale. Vi fokuserer på dyb vejrtrækning, smidighed og rolige øvelser for ryg, skuldre og nakke. Tempoet er langsomt, så både nybegyndere og stive kroppe kan følge med. Tilmeld dig på aftenro.dk.
\\
\midrule
\badge{NCDNavy}{E}\ \textbf{\color{NCDNavy} Lær at reparere din egen cykel}\newline
Bliver det dyrt at sende cyklen til mekaniker hver gang? På dette praktiske hold torsdag aften kl. 18.00--20.30 lærer du selv at holde din cykel kørende. Du medbringer din egen cykel, og alt værktøj findes på stedet. Vi gennemgår lapning af punkterede dæk, smøring af kæde samt justering af bremser og gear. Alle kan være med uanset erfaring. Skriv til vaerksted@bycykel.dk for tilmelding.
&
\badge{NCDNavy}{I}\ \textbf{\color{NCDNavy} Rejsespansk for begyndere -- Sprog \& Rejse}\newline
Skal du på ferie til Spanien eller Latinamerika, og vil du gerne kunne klare dig på det lokale sprog? På vores praktiske begynderhold om mandagen kl. 18.00--19.45 træner vi mundtlige hverdagssituationer som at købe billetter, bestille på restaurant og handle ind. Vi laver sjove rollespil i par og springer tunge grammatikbøger over. Der er ingen eksamen, kun god stemning! Kontakt os på sprogrejse.dk.
\\
\end{longtable}
\normalsize

% ---------------------------------------------------------------------
\subsubsection*{Opgave 2 (5 point)：Emma glemmer sin taske i S-toget}

\textbf{Instruktion:} Læs teksten \textit{Emma glemmer sin taske i S-toget}. I hvert afsnit er der én sætning, der ikke passer i afsnittet. Find den sætning, der ikke passer. Streg sætningen ud. Se eksemplet i afsnit (0).

\vspace{0.3em}
\textbf{\color{NCDNavy}Emma glemmer sin taske i S-toget}

\textbf{(0)} Fredag eftermiddag klokken 16.30 tog Emma S-toget fra Nørreport Station hjem mod Lyngby efter en lang dag på arbejde. Toget var næsten fyldt, så hun var glad for at finde en ledig siddeplads ved vinduet. Hun lagde sin grønne rygsæk op på hylden over sit sæde. \textcolor{NCDRed}{\textbf{[Eksempel — streges ud:]}}\ \textit{Emma kørte selv i sin bil hele vejen hjem.} Resten af turen lyttede hun til en podcast og kiggede ud ad vinduet.

\textbf{(1)} Da toget endelig standsede på Lyngby Station, rejste Emma sig hurtigt op og gik ud ad døren. Hun havde nemlig travlt, fordi hun skulle nå at handle før aftensmad. Først da hun stod ude på perronen, opdagede hun, at hun havde glemt sin rygsæk oppe på hylden i toget. Hun havde heldigvis sin rygsæk på ryggen hele tiden. Hun blev meget forskrækket, fordi både hendes pung, husnøgler og bærbare computer lå i tasken.

\textbf{(2)} Emma løb straks hen til billetkontoret på stationen for at få hjælp. En venlig DSB-medarbejder bad hende fortælle, hvilket tog hun var kommet med, og hvor hun havde siddet. Emma forklarede, at hun havde siddet i den forreste vogn på linje A mod Hillerød. Medarbejderen sagde, at der ikke kører S-tog til Hillerød. Derefter ringede han med det samme til personalet på endestationen i Hillerød for at efterlyse den grønne taske.

\textbf{(3)} Mens Emma ventede på svar fra Hillerød, gik hun nervøst frem og tilbage på stationen. Hun var især bekymret for sin computer, fordi hun havde mange vigtige dokumenter fra arbejdet på den. Emma var helt sikker på, at tasken var tom. Heldigvis ringede medarbejderens telefon allerede efter ti minutter. En flink passager havde nemlig fundet rygsækken på hylden og afleveret den til togføreren.

\textbf{(4)} Emma blev utrolig lettet og takkede DSB-medarbejderen mange gange. Han fortalte, at hun kunne hente sin rygsæk på kontoret i Hillerød hele aftenen indtil klokken 22.00. Kontoret i Hillerød var allerede lukket for dagen. Emma steg derfor på det næste S-tog nordpå og sad hele turen og kiggede spændt ud ad vinduet. Tyve minutter senere var hun fremme i Hillerød.

\textbf{(5)} På stationen i Hillerød viste Emma sit sundhedskort, og den unge mand bag skranken gav hende den grønne rygsæk tilbage. Hun åbnede lynlåsen med det samme for at tjekke, om alting stadig lå der. Desværre var både hendes computer og hendes pung væk. Alt var heldigvis på sin plads, så hun smilede stort og satte sig glad i toget hjem mod Lyngby. Fra den dag lægger Emma altid sin taske på gulvet mellem sine fødder.

% ---------------------------------------------------------------------
\subsubsection*{Opgave 3 (6 point)：Genbrug og loppemarkeder i Danmark}

\textbf{Instruktion:} Læs teksten \textit{Genbrug og loppemarkeder i Danmark}. Skriv de ord, der mangler. Du skal kun bruge hvert ord én gang. Der er fire ord, du ikke skal bruge. Se eksemplet i første linje.

\vspace{0.3em}
\textbf{\color{NCDNavy}Genbrug og loppemarkeder i Danmark}

\vspace{0.2em}
I Danmark \underline{\textbf{køber}} mange mennesker brugte ting i stedet for altid at vælge noget helt nyt. Det er både godt for miljøet og billigere end nye varer, og (1)\ \rule{2.1cm}{0.4pt} er genbrug blevet utrolig populært over hele landet. I næsten alle byer ligger der store genbrugsbutikker, (2)\ \rule{2.1cm}{0.4pt} frivillige sælger alt fra tøj og bøger til gamle lamper og møbler. Om sommeren går mange familier også på loppemarked i weekenden. Når vejret (3)\ \rule{2.1cm}{0.4pt} godt, stiller folk små borde op ude på gader og pladser. Her kan man finde sjove ting, (4)\ \rule{2.1cm}{0.4pt} ikke findes i almindelige butikker længere. Det handler dog (5)\ \rule{2.1cm}{0.4pt} kun om at spare penge. Mange danskere synes simpelthen, at gamle ting har mere personlighed, og de kan (6)\ \rule{2.1cm}{0.4pt} lide at give brugte møbler et nyt liv i hjemmet. På den måde mødes naboer over en kop kaffe og snakker sammen i solen.

\begin{grammatikbox}{Ordbank（选词框 · 共 11 词：1 例词 + 6 答案 + 4 多余干扰词）}
\centering
\large\bfseries\color{NCDNavy}
som \qquad havde \qquad godt \qquad derfor \qquad men \qquad køber \qquad hvor \qquad eller \qquad ikke \qquad hvis \qquad er
\end{grammatikbox}

% ---------------------------------------------------------------------
\subsubsection*{Opgave 4 (NY · 7 point)：Tre personer fortæller om deres transport til arbejde}

\textbf{Instruktion:}
\begin{itemize}
  \item Læs spørgsmålene (1--7), og læs de tre tekster (A, B og C).
  \item Find den person, der passer til hvert af de syv spørgsmål.
  \item Sæt kryds (\textbf{X}) ud for den person, der passer til spørgsmålet. Se eksemplet (0).
\end{itemize}

\vspace{0.2em}
\small
\renewcommand{\arraystretch}{1.3}
\begin{longtable}{@{} >{\raggedright\arraybackslash}p{11.3cm} @{\hspace{0.3cm}} >{\centering\arraybackslash}p{1.7cm} @{\hspace{0.15cm}} >{\centering\arraybackslash}p{1.7cm} @{\hspace{0.15cm}} >{\centering\arraybackslash}p{1.7cm} @{}}
\toprule
\textbf{Spørgsmål} & \textbf{A. Henrik} & \textbf{B. Nanna} & \textbf{C. Sami} \\
\midrule
\endfirsthead
\toprule
\textbf{Spørgsmål（续表）} & \textbf{A. Henrik} & \textbf{B. Nanna} & \textbf{C. Sami} \\
\midrule
\endhead
\bottomrule
\endfoot
0. Hvem havde en bil før i tiden, men har solgt den? \textcolor{NCDGray}{\small (Eksempel)} & \textbf{X} & B & C \\
\midrule
1. Hvem tager bussen i weekenden, selvom personen normalt kører i bil på arbejde? & A & B & C \\
\midrule
2. Hvem bruger tiden under transporten til at læse eller arbejde? & A & B & C \\
\midrule
3. Hvem kører sammen med en kollega for at dele udgifterne? & A & B & C \\
\midrule
4. Hvem har købt særligt tøj til dårligt vejr på grund af sin transport? & A & B & C \\
\midrule
5. Hvem oplever nogle gange forsinkelser på vejen til arbejde? & A & B & C \\
\midrule
6. Hvem har ikke noget kørekort? & A & B & C \\
\midrule
7. Hvem bor under ti kilometer fra sin arbejdsplads? & A & B & C \\
\end{longtable}
\normalsize

\vspace{0.3em}
\begin{center}
\textbf{\large\color{NCDNavy} Tre personer fortæller om deres transport til arbejde}
\end{center}

\begin{emailbox}{A. Henrik}
Jeg arbejder på et kontor midt i Aarhus, og jeg bor cirka ni kilometer uden for byen. Tidligere kørte jeg i bil hver dag, men jeg blev træt af at sidde i kø hver morgen og betale dyrt for parkering. Derfor solgte jeg min bil for et år siden og købte en god elcykel i stedet. Nu cykler jeg til arbejde året rundt, også når det regner eller blæser om vinteren. Jeg har bare købt varmt tøj og en vandtæt jakke. Turen tager kun femogtyve minutter, og jeg får både frisk luft og daglig motion. Desuden sparer jeg næsten to tusinde kroner om måneden. Jeg savner slet ikke min gamle bil, og jeg kommer faktisk hurtigere frem i myldretiden nu.
\end{emailbox}

\begin{emailbox}{B. Nanna}
Jeg bor i Roskilde og arbejder som folkeskolelærer i København, mere end tredive kilometer væk. Jeg har faktisk aldrig taget kørekort, så jeg tager regionaltoget til arbejde hver morgen. Togturen tager omkring femogtredive minutter, og bagefter går jeg ti minutter fra hovedbanegården hen til skolen. Jeg elsker at køre med tog, fordi jeg kan bruge tiden på at læse en god roman eller forberede mine timer til eleverne. Om efteråret er togene desværre tit forsinkede på grund af blade på skinnerne, og det kan være lidt irriterende. Alligevel vil jeg ikke bytte toget ud med en bil, fordi jeg slipper for at holde øje med trafikken. Min kæreste og jeg håber dog på at finde en lejlighed tættere på København næste år, så jeg kan gå på arbejde.
\end{emailbox}

\begin{emailbox}{C. Sami}
Jeg arbejder på et stort lager i et industriområde tyve kilometer fra min lejlighed i Odense. Jeg møder allerede klokken seks om morgenen, og på det tidspunkt kører der næsten ingen busser ud til industriområdet. Hvis jeg tager offentlig transport, skal jeg skifte bus to gange, og turen varer over en time. Derfor købte jeg en lille brugt bil for to år siden. Nu tager køreturen kun tyve minutter på motorvejen. Jeg har en kollega, Emil, som bor i samme gade som mig, så vi kører næsten altid sammen om morgenen. Han betaler halvdelen af benzinen, og vi hygger os med at høre radio på vejen. Når jeg skal ind til centrum i weekenden, tager jeg dog altid bussen, fordi parkering er besværlig.
\end{emailbox}

% =====================================================================
% DEL 2: SKRIVNING
% =====================================================================
\ncdsubsec{Del 2: Skrivning（书面表达｜minimum 70 ord）}

\subsubsection*{Opgave 1（官方必考 · 私人邮件回复）}

\textbf{Situation:} Du har fået en e-mail fra din ven, Sofie.\\
\textbf{Opgave:} Læs e-mailen, og skriv et svar til Sofie. Du skal svare på alle Sofies spørgsmål. Du skal skrive \textbf{minimum 70 ord}.

\begin{emailbox}{Modtaget e-mail fra Sofie}
Hej!

Jeg har hørt, at du har haft en god fødselsdagsfest for nylig! Tillykke med fødselsdagen! Hvordan gik det, og hvem kom? Hvad spiste og lavede I sammen? Og hvad fik du i gave --- eller hvad var det bedste ved aftenen?

Mange hilsner\\
Sofie
\end{emailbox}

\begin{emailbox}{Dit svar｜Til: Sofie \qquad Emne: Fødselsdag}
\vspace{0.3em}
\rule{\linewidth}{0.4pt}\\[0.55em]
\rule{\linewidth}{0.4pt}\\[0.55em]
\rule{\linewidth}{0.4pt}\\[0.55em]
\rule{\linewidth}{0.4pt}\\[0.55em]
\rule{\linewidth}{0.4pt}\\[0.55em]
\rule{\linewidth}{0.4pt}\\[0.55em]
\rule{\linewidth}{0.4pt}\\[0.55em]
\rule{\linewidth}{0.4pt}
\end{emailbox}

\subsubsection*{Bonus Opgave 1B（半正式实用邮件拓展｜礼貌请假与称谓一致性）}

\textbf{Situation:} Du går på et aftenkursus i madlavning (aftenskole) hver tirsdag aften. Næste tirsdag kan du desværre ikke komme, fordi du skal på en arbejdsrejse (eller besøge din familie).\\
\textbf{Opgave:} Skriv en høflig, halvt formel e-mail til din lærer, Mette. Du skal forklare, hvorfor du ikke kan komme, spørge om lektier eller materiale til næste gang, og bekræfte, hvornår du er tilbage igen. Husk at bruge konsekvent \textit{du / dig / din} over for din lærer. Skriv \textbf{minimum 70 ord}.

\begin{emailbox}{Dit svar｜Til: Mette \qquad Emne: Afbud til undervisningen på tirsdag}
\vspace{0.3em}
\rule{\linewidth}{0.4pt}\\[0.55em]
\rule{\linewidth}{0.4pt}\\[0.55em]
\rule{\linewidth}{0.4pt}\\[0.55em]
\rule{\linewidth}{0.4pt}\\[0.55em]
\rule{\linewidth}{0.4pt}\\[0.55em]
\rule{\linewidth}{0.4pt}\\[0.55em]
\rule{\linewidth}{0.4pt}
\end{emailbox}

% =====================================================================
% DEL 3: MUNDTLIG KOMMUNIKATION
% =====================================================================
\ncdsubsec{Del 3: Mundtlig kommunikation（口语交际全真题卡｜Opgave 1 \& 2）}

\subsubsection*{Mundtlig Opgave 1 (Emneopgave med Mindmap｜Del 1: 1--2 min. + Del 2: 3--4 min.)}

\begin{tipbox}{考场规则提醒（Opgave 1）}
\small
考生需自带\textbf{两张}提前填写好的思维导图（至少 1 张属于工作/实习/学习/语言学校领域）。导图框内\textbf{只能写关键词或短语（stikord），严禁写完整句子}！抽签选定 1 张后，先进行 1--2 分钟个人陈述（Del 1），随后回答考官 3--4 分钟追问（Del 2），两部分均须达到及格标准。
\end{tipbox}

\begin{grammatikbox}{Mindmap 1A（已填关键词卡）：Mit arbejde}
\small
\renewcommand{\arraystretch}{1.35}
\begin{tabularx}{\linewidth}{@{} >{\leavevmode\bfseries\color{NCDNavy}\raggedright\arraybackslash}p{3.2cm} @{\hspace{0.2cm}} >{\raggedright\arraybackslash}X @{\hspace{0.4cm}} >{\leavevmode\bfseries\color{NCDNavy}\raggedright\arraybackslash}p{3.0cm} @{\hspace{0.2cm}} >{\raggedright\arraybackslash}X @{}}
\toprule
\multicolumn{4}{c}{\large\bfseries\color{NCDRed} Emne: Mit arbejde \quad\small\color{NCDGray}（仅含 stikord 关键词，不含整句）} \\
\midrule
1. job / arbejdssted & kundeservice\-medarbejder · firma ”Nordic Lys“ · Aarhus C &
2. arbejdsopgaver & telefonsamtaler · e-mails til kunder · nye ordrer · problemløsning \\
\midrule
3. dage / tid & fuldtid · mandag--fredag · kl. 8.30--16.00 · fri i weekenden &
4. kollegaer / chef & 8 gode kollegaer · hjælpsom chef (Mads) · god stemning \\
\midrule
5. pauser & frokostpause 30 min. kl. 12.00 · kantine · eftermiddagskaffe &
6. transport til arbejde & bycykel · 15 minutter · cykelsti langs vandet \\
\bottomrule
\end{tabularx}
\end{grammatikbox}

\textbf{Del 2 考官追问模拟（Mit arbejde · 4 题）：}
\begin{enumerate}
  \item \textit{Du siger, at du taler med kunder i telefonen og på e-mail --- hvad er det sværeste eller mest spændende ved dine arbejdsopgaver?}
  \item \textit{Hvordan er atmosfæren på dit arbejde, og taler I mest dansk eller engelsk sammen i pauserne?}
  \item \textit{Hvad plejer I at spise i kantinen, og ses du nogle gange med dine kollegaer efter arbejde?}
  \item \textit{Hvordan kommer du på arbejde om vinteren, hvis det sner eller blæser rigtig meget?}
\end{enumerate}

\begin{grammatikbox}{Mindmap 1B（已填关键词卡）：Min bolig}
\small
\renewcommand{\arraystretch}{1.35}
\begin{tabularx}{\linewidth}{@{} >{\leavevmode\bfseries\color{NCDNavy}\raggedright\arraybackslash}p{3.2cm} @{\hspace{0.2cm}} >{\raggedright\arraybackslash}X @{\hspace{0.4cm}} >{\leavevmode\bfseries\color{NCDNavy}\raggedright\arraybackslash}p{3.0cm} @{\hspace{0.2cm}} >{\raggedright\arraybackslash}X @{}}
\toprule
\multicolumn{4}{c}{\large\bfseries\color{NCDRed} Emne: Min bolig \quad\small\color{NCDGray}（仅含 stikord 关键词，不含整句）} \\
\midrule
1. sted / by & Viby J · syd for Aarhus · tæt på letbanen &
2. hus / lejlighed & lys lejlighed · 2. sal · lille altan · gammel murstensejendom \\
\midrule
3. værelser / m\textsuperscript{2} & 3 værelser · 74 m\textsuperscript{2} · stue, soveværelse, køkken og bad &
4. alene / sammen med & sammen med kæreste (David) · indflytning sidste år \\
\midrule
5. naboer & flinke og rolige · ældre par ved siden af · hjælp til pakker &
6. område & grønt og trygt · lille park · supermarked og bager 200 m væk \\
\bottomrule
\end{tabularx}
\end{grammatikbox}

\textbf{Del 2 考官追问模拟（Min bolig · 4 题）：}
\begin{enumerate}
  \item \textit{Du nævner, at I har en lille altan --- bruger I den meget om sommeren, og hvad står der på den?}
  \item \textit{Hvordan fandt I lejligheden i Viby, og hvorfor valgte I netop det område?}
  \item \textit{Hvad kan du bedst lide ved dine naboer, og snakker I ofte sammen?}
  \item \textit{Vil I gerne blive boende i lejligheden i mange år, eller drømmer I om et hus en dag?}
\end{enumerate}

% ---------------------------------------------------------------------
\subsubsection*{Mundtlig Opgave 2 (Informationskløft｜满分 6 分，4--6 分及格)}

\begin{tipbox}{Opgave 2 考场核心守则（Case: Sofie）}
\small
本大题采取搭档配对形式，\textbf{仅对提问方（Kursist A）评分}：每提出 1 个聚焦第三人称问题得 1 分，共 6 格（4--6 分及格）。\\
\textbf{致命提醒}：题卡右上角的\textbf{人物头像照片（Portrætfotoet øverst）仅用于告知人物性别与姓名}，\textbf{切勿对头像本身提问}（如勿问 \textit{Hvor gammel er hun?}）！必须紧扣下方 2 个主题共 6 个图文格向搭档（Kursist B）提问，且双方严禁互相偷看题卡或提示对方如何发问。
\end{tipbox}

\begin{grammatikbox}{Kursist A Ark（提问方题卡 · 对准 6 个图文格发问，不可偷看 B 卡）}
\small
\textbf{Portrætfoto øverst:} Sofie（一位年轻女性微笑头像 — 仅作人物身份识别，\textbf{不要对头像提问}！）\par
\vspace{0.3em}
\renewcommand{\arraystretch}{1.35}
\begin{tabularx}{\linewidth}{@{} >{\leavevmode\bfseries\color{NCDNavy}\raggedright\arraybackslash}p{1.2cm} >{\leavevmode\bfseries\color{NCDTeal}\raggedright\arraybackslash}p{3.0cm} >{\raggedright\arraybackslash}X >{\centering\arraybackslash}p{1.0cm} @{}}
\toprule
\multicolumn{4}{@{}l}{\textbf{\color{NCDRed}Tema 1: Sofie arbejder i en børnehave}} \\
\midrule
\textbf{格 1} & \textit{sted / by} & 【图片提示】地图定位图标与幼儿园大门路牌 & \textbf{\large ?} \\
\textbf{格 2} & \textit{arbejdstider} & 【图片提示】工作日历（周一至周五）与时钟表盘 & \textbf{\large ?} \\
\textbf{格 3} & \textit{arbejdsopgaver} & 【图片提示】大人与小朋友一起读绘本、做手工和在户外滑梯旁玩耍 & \textbf{\large ?} \\
\midrule
\multicolumn{4}{@{}l}{\textbf{\color{NCDRed}Tema 2: Sofies fritid}} \\
\midrule
\textbf{格 4} & \textit{aktivitet} & 【图片提示】室内体育馆球拍与羽毛球图标 & \textbf{\large ?} \\
\textbf{格 5} & \textit{tid / dage} & 【图片提示】周历标出星期四及晚间时钟指针（19:00--20:30） & \textbf{\large ?} \\
\textbf{格 6} & \textit{alene / sammen med} & 【图片提示】几名好友并肩拿着球拍微笑的剪影 & \textbf{\large ?} \\
\bottomrule
\end{tabularx}
\end{grammatikbox}

\begin{emailbox}{Kursist B Ark（回答方信息卡 · 根据以下图文事实用第三人称作答）}
\small
\renewcommand{\arraystretch}{1.32}
\begin{tabularx}{\linewidth}{@{} >{\leavevmode\bfseries\color{NCDNavy}\raggedright\arraybackslash}p{4.8cm} >{\raggedright\arraybackslash}X @{}}
\toprule
\textbf{Tema / Kategori} & \textbf{B 卡实景图文信息（B 只能回答被问到的对应格）} \\
\midrule
\textbf{Tema 1.1: sted / by} & Odense（市中心路牌 + 幼儿园名牌：\textit{Børnehaven Solstrålen}） \\
\textbf{Tema 1.2: arbejdstider} & Mandag--fredag kl. 8.00--15.30（37 timer / fuldtid，周末双休） \\
\textbf{Tema 1.3: arbejdsopgaver} & Lege med børnene, læse bøger højt og gå tur på legepladsen \\
\midrule
\textbf{Tema 2.4: aktivitet} & Spiller badminton i en lokal sportsklub \\
\textbf{Tema 2.5: tid / dage} & Hver torsdag aften fra kl. 19.00 til 20.30 \\
\textbf{Tema 2.6: alene / sammen med} & Sammen med sin søster og to veninder \\
\bottomrule
\end{tabularx}
\end{emailbox}

% =====================================================================
% DEL 4: RETTEARK, DETALJERET ANALYSE & HØJSCORE-MODELLER
% =====================================================================
\ncdsubsec{Del 4: Retteark, Detaljeret Analyse \& Højscore-Modeller}

\subsubsection*{4.1 官方 Retteark 评分卡（Læsning 24 point \& Mundtlig Opgave 2）}

\small
\renewcommand{\arraystretch}{1.28}
\begin{longtable}{@{} >{\leavevmode\bfseries\color{NCDNavy}\raggedright\arraybackslash}p{3.2cm} @{\hspace{0.25cm}} >{\raggedright\arraybackslash}p{9.55cm} @{\hspace{0.25cm}} >{\centering\arraybackslash}p{1.8cm} @{\hspace{0.25cm}} >{\centering\arraybackslash}p{1.7cm} @{}}
\toprule
\textbf{模块 / 题号} & \textbf{Rettenøgle（标准答案）} & \textbf{分值} & \textbf{自测得分} \\
\midrule
\endfirsthead
\toprule
\textbf{模块 / 题号（续表）} & \textbf{Rettenøgle（标准答案）} & \textbf{分值} & \textbf{自测得分} \\
\midrule
\endhead
\bottomrule
\endfoot
\textbf{Læsning Opgave 1}\newline（6 题） &
1. \textbf{G} \quad 2. \textbf{C} \quad 3. \textbf{I} \quad 4. \textbf{E} \quad 5. \textbf{B} \quad 6. \textbf{H}\newline
\textcolor{NCDGray}{\small （多余干扰广告：D, F｜例题 0 为 A）} & 6 point & \rule{0.75cm}{0.4pt} / 6 \\
\midrule
\textbf{Læsning Opgave 2}\newline（5 题 · 须划掉） &
(1) \textit{Hun havde heldigvis sin rygsæk på ryggen hele tiden.}\newline
(2) \textit{Medarbejderen sagde, at der ikke kører S-tog til Hillerød.}\newline
(3) \textit{Emma var helt sikker på, at tasken var tom.}\newline
(4) \textit{Kontoret i Hillerød var allerede lukket for dagen.}\newline
(5) \textit{Desværre var både hendes computer og hendes pung væk.} & 5 point & \rule{0.75cm}{0.4pt} / 5 \\
\midrule
\textbf{Læsning Opgave 3}\newline（6 题） &
1. \textbf{derfor} \quad 2. \textbf{hvor} \quad 3. \textbf{er} \quad 4. \textbf{som} \quad 5. \textbf{ikke} \quad 6. \textbf{godt}\newline
\textcolor{NCDGray}{\small （多余干扰词：eller, havde, hvis, men｜例词 0 为 køber）} & 6 point & \rule{0.75cm}{0.4pt} / 6 \\
\midrule
\textbf{Læsning Opgave 4}\newline（7 题） &
1. \textbf{C} \quad 2. \textbf{B} \quad 3. \textbf{C} \quad 4. \textbf{A} \quad 5. \textbf{B} \quad 6. \textbf{B} \quad 7. \textbf{A}\newline
\textcolor{NCDGray}{\small （例题 0 为 A）} & 7 point & \rule{0.75cm}{0.4pt} / 7 \\
\midrule
\textbf{Læsning 总计} & \textbf{Bestået：16--24 point（及格）} \quad|\quad \textbf{Ikke bestået：0--15 point} & \textbf{24 point} & \rule{0.75cm}{0.4pt} / 24 \\
\midrule
\textbf{Mundtlig Opg. 2} & 6 个聚焦第三人称问题，每格 1 分（\textbf{Bestået：4--6 point}） & \textbf{6 point} & \rule{0.75cm}{0.4pt} / 6 \\
\end{longtable}
\normalsize

\begin{tipbox}{Skrivning \& Mundtlig Opgave 1 自检评分准则}
\small
\begin{itemize}
  \item \textbf{Skrivning（及格底线）}：\ding{172} 字数严格达到 \textbf{70 词以上}（建议写到 85--110 词，避免因缩写连词扣字数）；\ding{173} \textbf{完整回答全部 4 个问题}（不能漏问）；\ding{174} 时态一致（回忆生日用过去时 \textit{var / spiste / fik}，请假说明用现在时与将来表达 \textit{kan ikke / skal / kommer}）；\ding{175} 从句 \textit{ikke} 置于动词前，句首状语触发 V2 倒装；\ding{176} 邮件称呼与落款格式地道。
  \item \textbf{Mundtlig Opgave 1}：Del 1 围绕 6 个 \textit{stikord} 流畅陈述 1--2 分钟，Del 2 能自然接住考官 3--4 分钟的追问（哪怕短答后主动加一句理由/例子），两部分均须及格。
\end{itemize}
\end{tipbox}

% ---------------------------------------------------------------------
\subsubsection*{4.2 阅读 Opgave 1--4 逐题同义转述解析、全文精译与核心词表}

\paragraph*{一、Opgave 1（成人夜校与业余课程）同义转述定位与干扰项拆解}
\begin{itemize}
  \item \textbf{(0) Jonas $\rightarrow$ A [Eksempel]}：\textit{lave noget kreativt med sine hænder} $\leftrightarrow$ \textit{arbejde med hænderne}；\textit{kopper og skåle i ler} $\leftrightarrow$ \textit{kopper, tallerkener og skåle i blødt ler}；\textit{kun tid i weekenden} $\leftrightarrow$ \textit{om lørdagen kl. 10.00--13.00}；\textit{aldrig prøvet det før} $\leftrightarrow$ \textit{ingen tidligere erfaring}。
  \item \textbf{1. Astrid $\rightarrow$ G}：同义定位：\textit{ikke et stort kamera} $\leftrightarrow$ \textit{uden et dyrt spejlreflekskamera}；\textit{kort aftenhold for begyndere} $\leftrightarrow$ \textit{korte kursus over fire tirsdag aftener ... for alle uden erfaring}；\textit{bruge telefonens kamera og redigere fotos} $\leftrightarrow$ \textit{feriebilleder på din smartphone ... hurtig billedredigering}。\textbf{排除干扰项 D}：D 属于同主题摄影陷阱，但明确要求自带“带可换镜头和三脚架的单反相机（\textit{spejlreflekskamera / systemkamera}）”，且为周末白天长达 12 周的进阶户外课，与 Astrid 的条件完全相反。
  \item \textbf{2. Tobias $\rightarrow$ C}：同义定位：\textit{spise mindre kød} $\leftrightarrow$ \textit{retter uden kød}；\textit{varme grøntsagsretter med gode krydderier} $\leftrightarrow$ \textit{varme retter med bønner, linser, friske grøntsager og aromatiske krydderier}；\textit{spiser aftensmaden bagefter} $\leftrightarrow$ \textit{slutter altid aftenen med at spise ved et langt bord sammen}；\textit{onsdagen eller torsdagen efter arbejde} $\leftrightarrow$ \textit{hver onsdag kl. 17.30--20.30}。\textbf{排除干扰项 F}：F 同为厨房/美食领域干扰项，但 F 教做面包和蛋糕、上课在周日上午且将成品打包带回家，并不满足无肉热菜晚餐、聚餐与周三/四晚间时间。
  \item \textbf{3. Yasmin $\rightarrow$ I}：同义定位：\textit{rejse ... i Spanien} $\leftrightarrow$ \textit{på ferie til Spanien eller Latinamerika}；\textit{ikke interesseret i svær grammatik eller skriftlige prøver} $\leftrightarrow$ \textit{springer tunge grammatikbøger over ... ingen eksamen}；\textit{snakke, bestille mad på café og spørge om vej} $\leftrightarrow$ \textit{mundtlige hverdagssituationer som at købe billetter, bestille på restaurant og handle ind}。其余广告无语言课程。
  \item \textbf{4. Anders $\rightarrow$ E}：同义定位：\textit{alt for mange penge hos cykelsmeden} $\leftrightarrow$ \textit{dyrt at sende cyklen til mekaniker}；\textit{tage sin egen cykel med} $\leftrightarrow$ \textit{medbringer din egen cykel}；\textit{lappe slanger, rense kæden og justere gear} $\leftrightarrow$ \textit{lapning af punkterede dæk, smøring af kæde samt justering af bremser og gear}。
  \item \textbf{5. Birgitte $\rightarrow$ B}：同义定位：\textit{kan desværre ikke læse noder} $\leftrightarrow$ \textit{tror du, at man skal kunne læse noder? ... uden noder}；\textit{ikke til en svær optagelsesprøve eller synge alene} $\leftrightarrow$ \textit{ingen synger solo, og der er ingen optagelsesprøve}；\textit{kendte danske sange og popsange} $\leftrightarrow$ \textit{kendte danske klassikere og nyere pophits}。
  \item \textbf{6. Omar $\rightarrow$ H}：同义定位：\textit{ondt i nakken og ryggen} $\leftrightarrow$ \textit{øvelser for ryg, skuldre og nakke}；\textit{ønsker ikke hård fitness med høj musik eller konkurrence} $\leftrightarrow$ \textit{komme ned i tempo ... stille, varmt lokale ... tempoet er langsomt}；\textit{blide strækøvelser, træner vejrtrækning} $\leftrightarrow$ \textit{blid yoga ... dyb vejrtrækning, smidighed}。
\end{itemize}

\begin{transbox}
\small
\textbf{\color{NCDTeal}[Opgave 1 人物要求与广告中文精译]}\\
\textbf{人物 (0)--(6)}：\textbf{(0) 约纳斯}想用双手做些创意手工，但周一至周五每天加班到很晚，只有周末有空，梦想从零学习用陶土捏杯子和碗。\textbf{1. 阿斯特丽德}（62 岁）刚买了新手机，想拍更清晰的孙辈和度假照片，没有大相机也不想报长课，想找个零基础短晚班学习手机拍摄和修图。\textbf{2. 托比亚斯}喜欢下厨但想平时少吃肉，寻找带香料的蔬菜热菜新食谱，看重团队一起做完后围坐共进晚餐，且只有周三或周四下班后有空。\textbf{3. 亚斯敏}夏天要去西班牙旅行三周，想学些日常口语，对枯燥语法和笔试没兴趣，想找练点餐问路的零基础口语班。\textbf{4. 安德斯}每天骑车 12 公里上班，修车花钱太多，想找晚间实操工坊带自己的车去学补胎、洗链条和调变速。\textbf{5. 比吉特}热爱音乐但不会看乐谱，也不想参加严格的独唱考核，想在轻松的晚间合唱团唱丹麦金曲和流行歌。\textbf{6. 奥马尔}每天对着电脑 8 小时导致肩颈背酸痛，不喜欢嘈杂激烈的健身房，想在晚上参加舒缓拉伸、练呼吸、让身心平静下来的课程。\\[3pt]
\textbf{广告 A--I}：\textbf{A（例）周末陶艺}：奥胡斯陶艺工坊周六 10:00--13:00，零基础学捏杯盘碗及上釉，含材料费与烧制费。\textbf{B 大家都能唱！无谱流行合唱团}：Klang \& Glæde 每周二 19:00--21:00，无需识谱、无独唱、无入学试音，唱丹麦经典与流行热曲，中场休息一起喝咖啡。\textbf{C 绿色日常餐}：中城区烹饪学校每周三 17:30--20:30，烹饪豆类、扁豆、鲜蔬与香料无肉热菜，分组合作后长桌共进晚餐，含食材费。\textbf{D 单反自然摄影（干扰项）}：东日德兰摄影俱乐部 12 周课程，每周六上午 9:00--13:00 赴林海拍鸟类风景，须自带可换镜头单反与三脚架并具备手动曝光基础。\textbf{E 学修自己的自行车}：周四晚 18:00--20:30，自带自行车，现场提供全部工具，零基础学补胎、上链条油、调刹车和变速器。\textbf{F 酸面包与法式蛋糕（干扰项）}：港口街烘焙坊每周日上午 10:00--14:00，教揉面与调温巧克力，全部烤好的面包蛋糕打包带回家，四周 950 克朗。\textbf{G 手机摄影}：文化馆夜校四周短班，每周二晚 18:30--20:30，无需昂贵单反，零基础学习手机光线构图、人像度假照及免费修图软件。\textbf{H 舒缓瑜伽与拉伸}：Aftenro 运动中心周一、周三晚 19:30--20:45，安静温暖教室，慢节奏深呼吸、提升柔韧性并缓解背肩颈僵硬。\textbf{I 零基础旅游西语}：周一晚 18:00--19:45，情景对话练习买票、餐厅点菜与购物，跳过繁重语法书、无考试。
\end{transbox}

\paragraph*{二、Opgave 2（城铁遗落背包记）矛盾句定位与全文精译}
\begin{itemize}
  \item \textbf{(0) [Eksempel]}：开头已说明艾玛乘坐 S-tog 城铁，因此“自己开车一路回家”与交通方式矛盾。
  \item \textbf{(1)}：\textit{Hun havde heldigvis sin rygsæk på ryggen hele tiden.} —— 前一句刚说她在站台发现“把背包忘在列车行李架上了”，后一句又紧接“紧张钱包电脑都在包里”，与“包一直背在背上”直接矛盾。
  \item \textbf{(2)}：\textit{Medarbejderen sagde, at der ikke kører S-tog til Hillerød.} —— 前一句说明她坐的是开往 Hillerød 的 A 线，下一句工作人员也立刻打电话给终点站 Hillerød 寻包，因此“没有开往 Hillerød 的城铁”矛盾。
  \item \textbf{(3)}：\textit{Emma var helt sikker på, at tasken var tom.} —— 前一句说她最担心包里装有重要工作文件的电脑，且第 (1) 段交代包里有钱包、钥匙和电脑，故“完全确定包是空的”矛盾。
  \item \textbf{(4)}：\textit{Kontoret i Hillerød var allerede lukket for dagen.} —— 前一句刚告知“整个晚上一直到 22 点都可以去 Hillerød 办公室取包”，下一句她也立刻搭下班车北上取包，故“当天已经关门”矛盾。
  \item \textbf{(5)}：\textit{Desværre var både hendes computer og hendes pung væk.} —— 紧接着下一句是“幸好所有东西都在原位，她开心地笑了”，与“电脑和钱包都不见了”完全矛盾。
\end{itemize}

\begin{transbox}
\small
\textbf{\color{NCDTeal}[Opgave 2 全文中文精译（已剔除矛盾句）]}\\
\textbf{(0)} 周五下午 16:30，漫长的工作日结束后，艾玛从北门站（Nørreport）乘坐 S-tog 城铁回灵比（Lyngby）。车厢里几乎坐满了疲惫的通勤族，她很高兴找到一个靠窗的空位，并把绿色双肩包放上了座位上方的行李架。余下的车程她边听播客边看窗外。\textbf{(1)} 火车终于停在灵比站，艾玛迅速起身走出了车门，因为她急着回家前先去买点东西。直到站在站台上，她才猛然发现自己把背包忘在了车厢行李架上。她一下子吓坏了，因为钱包、家门钥匙和笔记本电脑全在包里。\textbf{(2)} 艾玛立刻跑向站台旁的售票办公室求助。一位友善的 DSB 工作人员请她说明搭乘的车次和座位位置。艾玛解释自己坐在开往希勒勒（Hillerød）方向的 A 线第一节车厢里。工作人员随即给终点站希勒勒的同事打电话协寻这个绿背包。\textbf{(3)} 在等待希勒勒消息期间，艾玛焦急地在车站来回踱步。她尤其担心那台电脑，里面存着许多重要的工作文件。幸运的是，短短十分钟后工作人员的电话就响了——原来一位好心乘客在架子上发现了背包并交给了列车长。\textbf{(4)} 艾玛如释重负，连连向 DSB 工作人员道谢。对方告诉她，今晚 22:00 前随时都可以去希勒勒车站办公室认领背包。于是艾玛登上了下一班北上的城铁，一路兴奋地看着窗外，二十分钟后便抵达了希勒勒。\textbf{(5)} 在希勒勒车站，艾玛出示了医保黄卡，柜台后的年轻小伙把绿色背包还给了她。她立刻拉开拉链检查东西是否都在。幸好每样物品都完好无损，她露出灿烂的笑容，开开心心地坐上了返回灵比的火车。从那天起，艾玛总是把包放在两脚之间的地板上。
\end{transbox}

\paragraph*{三、Opgave 3（丹麦二手与跳蚤市场文化）完形语法解析与全文精译}
\begin{itemize}
  \item \textbf{(1) derfor}：连接两个完整分句并引出结果（“既环保又省钱，\textbf{因此}二手循环在全国变得极受欢迎”）；后面紧跟倒装语序 \textit{er genbrug blevet}，符合副词置前触发 V2 倒装规则。干扰词 \textit{men}（转折）与前文因果不符；\textit{eller} / \textit{hvis} / \textit{havde} 均不合句法。
  \item \textbf{(2) hvor}：引导定语从句修饰地点名词 \textit{genbrugsbutikker}（“在那里志愿者出售……”），因从句主语（\textit{frivillige}）和宾语（\textit{alt...}）均已完整，故必须用关系副词 \textit{hvor}，不能用 \textit{som}。
  \item \textbf{(3) er}：在时间/条件从句中与表语 \textit{godt} 搭配，且主句动词 \textit{stiller} 为一般现在时，故只能用现在时系动词 \textit{er}（不可用过去时助动词 \textit{havde}，且丹麦语不说 *\textit{vejret har godt}）。
  \item \textbf{(4) som}：关系代词作从句主语，指代前面的物（\textit{sjove ting}），后接被动形式动词 \textit{ikke findes}（注意从句否定副词 \textit{ikke} 紧跟在 \textit{som} 后、动词前！）。
  \item \textbf{(5) ikke}：构成经典让步搭配 \textit{handler dog ikke kun om...}（“然而这并不仅仅关乎省钱”），承接下句“许多丹麦人纯粹觉得旧物更有个性”。
  \item \textbf{(6) godt}：高频固定搭配 \textit{kan godt lide at...}（“喜欢做某事”），不能用其他词替代。
\end{itemize}

\begin{transbox}
\small
\textbf{\color{NCDTeal}[Opgave 3 全文中文精译]}\\
在丹麦，许多人会购买二手物品，而不是总是挑选全新的东西。这不仅有益于环境，也比买新品更便宜，\textbf{因此}二手回收文化在全国各地变得非常流行。几乎每个城市都有大型二手公益店，\textbf{在那里}志愿者出售从衣物、书籍到老台灯和家具等各种商品。到了夏天，许多家庭还会在周末逛跳蚤市场。当天气好的时候，人们会在街道和广场上摆起小摊桌。在这里你能淘到普通商店里再也买不到的有趣物件。不过，这并\textbf{不}仅仅是为了省钱。许多丹麦人纯粹觉得老物件更有个性，而且他们非常喜欢让二手家具在家里焕发新生。通过这种方式，邻里们还能在阳光下喝着咖啡聊天相聚。
\end{transbox}

\paragraph*{四、Opgave 4（三人谈通勤方式）逐题定位解析与三人自述精译}
\begin{itemize}
  \item \textbf{(0) A [Eksempel]}：Henrik 明确说 \textit{Tidligere kørte jeg i bil... Derfor solgte jeg min bil for et år siden}（过去开车，一年前卖了车）。Nanna 从未考驾照；Sami 仍在开车上班。
  \item \textbf{1. C (Sami)}：平时开二手小车上班，但周末去市中心总是坐公交（\textit{Når jeg skal ind til centrum i weekenden, tager jeg dog altid bussen}）。
  \item \textbf{2. B (Nanna)}：在火车上利用通勤时间读小说或备课（\textit{bruge tiden på at læse en god roman eller forberede mine timer}），对应题干 \textit{læse eller arbejde}。
  \item \textbf{3. C (Sami)}：同街同事 Emil 搭顺风车并分摊一半油费（\textit{kører næsten altid sammen ... Han betaler halvdelen af benzinen}），对应 \textit{dele udgifterne}。
  \item \textbf{4. A (Henrik)}：因为全年雨雪天气骑车上班，专门买了保暖衣物和防水外套（\textit{købt varmt tøj og en vandtæt jakke}）。
  \item \textbf{5. B (Nanna)}：秋天轨道落叶常导致火车晚点（\textit{tit forsinkede}），对应题干 \textit{forsinkelser}。
  \item \textbf{6. B (Nanna)}：从未考过驾照（\textit{aldrig taget kørekort}）。Henrik 曾长期开车，Sami 每天开车。
  \item \textbf{7. A (Henrik)}：住在市中心外约 9 公里处（\textit{cirka ni kilometer} $< 10$ 公里），而 Nanna 超 30 公里，Sami 20 公里。
\end{itemize}

\begin{transbox}
\small
\textbf{\color{NCDTeal}[Opgave 4 三人自述中文精译]}\\
\textbf{A. 亨里克}：我在奥胡斯市中心的一家办公室工作，住在城外大约九公里处。以前我每天开车上班，但渐渐厌倦了每天早上堵在车流里还要付高昂的停车费。因此一年前我卖掉了汽车，买了一辆好用的电动自行车。现在我全年骑车上班，哪怕冬天刮风下雨也照骑不误。我只是添置了保暖衣物和一件防水外套。单程只需二十五分钟，既能呼吸新鲜空气，又能天天锻炼。此外，我每月还能省下将近两千克朗。我一点也不想念我的旧汽车，事实上现在高峰期我比以前到得还快。\\
\textbf{B. 南娜}：我住在罗斯基勒，在三十多公里外的哥本哈根担任公立学校教师。我其实从未考过驾照，所以每天早上搭区际火车上班。火车车程大约三十五分钟，下车后我再从中央车站步行十分钟到学校。我非常喜欢坐火车，因为我可以利用这段时间读一本好小说，或者为学生备课。可惜到了秋天，因为铁轨上有落叶，火车经常晚点，这有时让人有点烦躁。尽管如此，我也不想把火车换成汽车，因为坐火车不用盯着路况看。不过我和男朋友希望明年能在离哥本哈根更近的地方找到公寓，这样我就能步行上班了。\\
\textbf{C. 萨米}：我在欧登塞公寓二十公里外一个工业区的大型仓库工作。我早上六点就要到岗，那个时间点几乎没有公交车开往工业区。如果坐公共交通，我得换乘两次公交，单程超过一小时。因此两年前我买了一辆二手小汽车。现在走高速公路只要二十分钟。我有个住在同一条街上的同事埃米尔，所以我们早上几乎总是一起拼车上班。他付一半的油钱，我们在路上听着收音机聊得很开心。不过周末要进市中心的时候，因为停车很麻烦，我总是改坐公交车。
\end{transbox}

\paragraph*{五、核心高频词表（Gloser｜15 个必背屈折变化词条）}

\begin{glosertable}
名词 & en rygsæk & \textit{(rygsækken, rygsække, rygsækkene)} & 背包，双肩包 \\
名词 & en perron & \textit{(perronen, perroner, perronerne)} & 站台，月台 \\
动词 & at opdage & \textit{(opdager, opdagede, har opdaget)} & 发现，察觉 \\
动词 & at efterlyse & \textit{(efterlyser, efterlyste, har efterlyst)} & 报失协寻，通告寻找 \\
名词 & et loppemarked & \textit{(loppemarkedet, loppemarkeder, -ne)} & 跳蚤市场 \\
名词 & en genbrugsbutik & \textit{(genbrugsbutikken, -butikker, -kerne)} & 二手公益商店 \\
形容词 & frivillig & \textit{(frivillig, frivilligt, frivillige)} & 自愿的；志愿者（常用复数） \\
副词 & derfor / desuden / alligevel & （句首触发 V2；从句中在动词前） & 因此 / 此外 / 尽管如此 \\
动词 & at sælge & \textit{(sælger, solgte, har solgt)} & 出售，卖（不规则！） \\
名词 & et kørekort & \textit{(kørekortet, kørekort, kørekortene)} & 驾驶证（单复数同形） \\
形容词 & forsinket & \textit{(forsinket, forsinket, forsinkede)} & 晚点的，延误的 \\
名词 & en myldretid & \textit{(myldretiden, ingen flertal)} & 上下班高峰期 \\
动词 & at forberede & \textit{(forbereder, forberedte, har forberedt)} & 准备（\textit{forberede sig} 做准备） \\
动词 & at deltage & \textit{(deltager, deltog, har deltaget)} & 参加（搭配介词 \textit{i}） \\
名词 & en fødselsdagsgave & \textit{(fødselsdagsgaven, -gaver, -gaverne)} & 生日礼物 \\
\end{glosertable}

% =====================================================================
% 4.3 写作范文
% =====================================================================
\subsubsection*{4.3 写作双档高分示范与点评（Skrivning Modeller）}

\begin{emailbox}{Opgave 1 范文一：A2 稳过版（88 词｜清晰覆盖 4 个问题 · 时态准确）}
Hej Sofie

Tak for din søde mail! Ja, jeg havde en rigtig dejlig fødselsdagsfest i lørdags. Mine to bedste venner, min søster og tre kollegaer fra mit arbejde kom hjem til mig om aftenen. Vi spiste lækker italiensk pasta med salat, og bagefter drak vi kaffe og spiste chokoladekage. Vi snakkede meget, hørte god musik og spillede et sjovt brætspil sammen. Jeg fik en ny varm vinterjakke og nogle gode bøger i gave, men det bedste ved aftenen var, at alle hyggede sig så meget.

Mange hilsner\\
Anna

\vspace{0.3em}
\rule{\linewidth}{0.4pt}\\[0.2em]
\small\textbf{\color{NCDTeal}[考官点评 · 88 ord]}：完全满足 70 词底线，自然应答全部四个问题（聚会好不好、谁来了、吃喝和活动、礼物与最棒的瞬间）。全文一般过去时切换稳定（\textit{havde, kom, spiste, drak, snakkede, hørte, spillede, fik, var, hyggede sig}），句首状语 \textit{bagefter drak vi} 正确使用 V2 倒装，末尾包含规范的 \textit{at} 宾语从句。
\end{emailbox}

\begin{emailbox}{Opgave 1 范文二：A2+ 高分示范版（113 词｜多从句切换 + 过去完成时 + 否定前置）}
Hej Sofie

Tusind tak for din besked og din søde hilsen! Ja, jeg holdt en fantastisk fødselsdagsmiddag i lørdags, hvor jeg havde inviteret otte venner og min bror hjem. \da{Selvom vi ikke har set} hinanden så tit på det sidste, \dabold{blev stemningen} hurtigt varm og festlig. Vi startede med hjemmelavet lasagne og salat, og bagefter spiste vi en stor hindbærkage, \da{som} min veninde selv \dabold{havde bagt} til mig. Efter maden snakkede vi længe og spillede et sjovt kortspil, hvor alle grinede hele tiden. Jeg fik både en ny cykelhjelm og et gavekort til biografen, men det allerbedste ved aftenen var helt sikkert at fejre dagen sammen med så gode mennesker.

Kærlig hilsen\\
Anna

\vspace{0.3em}
\rule{\linewidth}{0.4pt}\\[0.2em]
\small\textbf{\color{NCDNavy}[高分亮点 · 113 ord]}：\ding{172} 句首让步从句 \textit{Selvom vi \textbf{ikke} har set...} 严格做到 \textit{ikke} 在助动词前，主句 \textit{blev stemningen} 紧跟 V2 倒装；\ding{173} 灵活穿插过去完成时 \textit{havde inviteret / havde bagt} 与一般过去时 \textit{holdt / spiste / fik} 形成背景层次；\ding{174} 关系从句 \textit{hvor / som} 衔接紧凑自然，完全具备冲刺 Modul 3 的文本组织力。
\end{emailbox}

\begin{emailbox}{Bonus Opgave 1B 范文：半正式夜校请假邮件（95 词｜du/dig 一致性与礼貌请求）}
Kære Mette

Jeg skriver til \da{dig}, fordi jeg desværre \dabold{ikke kan komme} til aftenkurset på tirsdag. Jeg skal nemlig på en kort arbejdsrejse til Hamborg fra mandag morgen til onsdag aften, så jeg når ikke tilbage til Aarhus i tide. Jeg er rigtig ked af at gå glip af undervisningen, da jeg er meget glad for holdet. \da{Vil du eventuelt sende} de vigtigste opskrifter eller noter til mig på mail, så jeg kan læse dem derhjemme? Jeg kommer helt sikkert igen tirsdag i næste uge og glæder mig allerede til næste gang.

Venlig hilsen\\
Anna

\vspace{0.3em}
\rule{\linewidth}{0.4pt}\\[0.2em]
\small\textbf{\color{NCDGold}[半正式语域要点 · 95 ord]}：对丹麦老师直接用名字和 \textit{Kære Mette / Venlig hilsen}，全程保持单数 \textit{dig / du} 称谓一致（切忌混用 \textit{De} 或 \textit{I}）。从句 \textit{fordi jeg desværre ikke kan komme} 中副词 \textit{desværre} 和 \textit{ikke} 双双前置在动词前；请求用 \textit{Vil du eventuelt...} 礼貌不生硬，且清晰交代了请假原因、求发资料与返校日期。
\end{emailbox}

% =====================================================================
% 4.4 口语示范稿
% =====================================================================
\subsubsection*{4.4 口语示范逐字稿、核心 IPA 与信息差双轨句型表}

\begin{transbox}
\small
\textbf{\color{NCDNavy}[Mindmap 1A 口语陈述范文：Mit arbejde｜138 词 · 约 1.5 分钟]}\\
\textit{Jeg vil gerne fortælle lidt om mit arbejde \ipa{ˈɑːbɑi̯ˀdə}. Jeg arbejder som kundeservicemedarbejder i et mellemstort firma i Aarhus centrum, som hedder Nordic Lys. Det er et fuldtidsjob \ipa{ˈfultˢiðsˌjʌb}, så jeg møder mandag til fredag fra klokken halv ni til klokken fire, og heldigvis har jeg altid fri i weekenden \ipa{ˈviːɡˌɛnˀdən}. Mine vigtigste arbejdsopgaver \ipa{ˈɑːbɑi̯ˀdsˌʌbˌɡæːvɐ} er at tale med kunder i telefonen, svare på e-mails og registrere nye ordrer i systemet. Når et produkt ikke virker, prøver jeg at finde en god løsning. Jeg har otte rigtig søde kollegaer \ipa{kʰoˈleːɐ} og en hjælpsom chef, som hedder Mads. I vores frokostpause \ipa{ˈfʁɔɡʌsdˌpʰɑu̯sə} klokken tolv spiser vi næsten altid sammen i kantinen, og om eftermiddagen drikker vi tit en kop kaffe. Jeg cykler \ipa{ˈsyɡlɐ} på arbejde hver dag langs vandet \ipa{ˈvanˀəð}, og turen tager kun femten minutter. Jeg er rigtig glad for mit arbejde.}
\end{transbox}

\paragraph*{Mindmap 1A Del 2 考官追问实战示范对话：}
\dialogturn{1}{Eksaminator}{Du siger, at du taler med kunder i telefonen og på e-mail — hvad er det sværeste ved dine arbejdsopgaver?}{sige \ipa{ˈsiːə} · sværeste \ipa{ˈsvɛːɐsdə}}{你说你通过电话和邮件与客户沟通——你工作任务中最难的是什么？}
\dialogturn{2}{Kursist}{Det sværeste er nogle gange at forstå kunderne hurtigt på dansk i telefonen, fordi man ikke kan se hinanden. Men hvis jeg spørger høfligt, om de vil tale lidt langsommere, er alle meget søde.}{telefonen \ipa{tˢeləˈfoːˀnən} · langsommere \ipa{ˈlɑŋˌsʌmˀɐ}}{有时最难的是在电话里快速听懂丹麦语，因为看不到对方。但如果我礼貌地请他们说慢一点，大家都会非常友善。}

\begin{transbox}
\small
\textbf{\color{NCDNavy}[Mindmap 1B 口语陈述范文：Min bolig｜141 词 · 约 1.5 分钟]}\\
\textit{Nu vil jeg gerne fortælle lidt om min bolig \ipa{ˈboːli}. Jeg bor i Viby, som er en hyggelig bydel syd for Aarhus, kun få minutter fra letbanen \ipa{ˈlɛdˌbæːnən}. Min kæreste David og jeg bor i en lys lejlighed \ipa{ˈlɑi̯liˌheðˀ} på anden sal i en gammel murstensejendom, og vi har også en lille altan \ipa{alˈtˢæːˀn}. Lejligheden er på fireoghalvfjerds kvadratmeter og har tre værelser: en stor stue, et soveværelse \ipa{ˈsʌu̯əˌvɛːˀəlsə}, et fint køkken og et lille badeværelse. Vi flyttede ind sammen sidste år, fordi vi gerne ville have mere plads. Vores naboer \ipa{ˈnæːˌboˀɐ} er meget flinke og rolige. Ved siden af os bor der et ældre par, som ofte tager imod vores pakker, når vi ikke er hjemme. Selve området \ipa{ˈʌmˌʁɔːˀəð} er grønt og trygt, og der ligger både en park, et supermarked og en bager to hundrede meter væk. Vi trives rigtig godt her.}
\end{transbox}

\paragraph*{Mindmap 1B Del 2 考官追问实战示范对话：}
\dialogturn{1}{Eksaminator}{Du nævner, at I har en lille altan — bruger I den meget om sommeren?}{nævner \ipa{ˈnɛu̯nɐ} · altan \ipa{alˈtˢæːˀn}}{你提到你们有个小阳台——夏天的时候你们常用吗？}
\dialogturn{2}{Kursist}{Ja, vi bruger den næsten hver aften, når vejret er godt! Vi har to små stole, et lille bord og nogle grønne planter med tomater, og vi drikker tit aftenskaffe derude.}{vejret \ipa{ˈvɛːˀɐð} · tomater \ipa{tˢoˈmæːˀdɐ}}{是的，只要天气好，我们几乎每天晚上都用！上面放着两把小椅子、一张小桌和几盆绿植西红柿，我们常在外面喝晚茶咖啡。}

\paragraph*{Mundtlig Opgave 2（Sofie）6 格提问双轨对照表与第三人称转述示范：}

\small
\renewcommand{\arraystretch}{1.28}
\begin{longtable}{@{} >{\leavevmode\bfseries\color{NCDNavy}\raggedright\arraybackslash}p{2.55cm} @{\hspace{0.25cm}} >{\raggedright\arraybackslash}p{4.0cm} @{\hspace{0.25cm}} >{\raggedright\arraybackslash}p{5.05cm} @{\hspace{0.25cm}} >{\raggedright\arraybackslash}p{4.65cm} @{}}
\toprule
\textbf{题格关键词} & \textbf{直接问法（稳拿满分）} & \textbf{礼貌间接问法（注意不倒装！）} & \textbf{Kursist B 标准作答} \\
\midrule
\endfirsthead
\toprule
\textbf{题格关键词（续表）} & \textbf{直接问法（稳拿满分）} & \textbf{礼貌间接问法（注意不倒装！）} & \textbf{Kursist B 标准作答} \\
\midrule
\endhead
\bottomrule
\endfoot
1. sted / by & \textit{I hvilken by arbejder Sofie, og hvad hedder børnehaven?} & \textit{Kan du fortælle mig, \da{hvor} \dabold{Sofie arbejder}?} & \textit{Hun arbejder i Børnehaven Solstrålen i Odense.} \\
\midrule
2. arbejdstider & \textit{Hvornår arbejder hun, og har hun fri i weekenden?} & \textit{Jeg vil gerne vide, \da{hvornår} \dabold{hun møder}, og \da{om} \dabold{hun arbejder} i weekenden.} & \textit{Hun arbejder mandag til fredag kl. 8.00--15.30 og har fri i weekenden.} \\
\midrule
3. arbejdsopgaver & \textit{Hvad er hendes arbejdsopgaver i børnehaven?} & \textit{Ved du, \da{hvad} \dabold{hun laver} sammen med børnene?} & \textit{Hun leger med børnene, læser bøger højt og går tur på legepladsen.} \\
\midrule
4. aktivitet & \textit{Hvad laver Sofie i sin fritid?} & \textit{Kan du sige mig, \da{hvilken sport} \dabold{hun dyrker} i sin fritid?} & \textit{Hun spiller badminton i en lokal sportsklub.} \\
\midrule
5. tid / dage & \textit{Hvilke dage og hvornår spiller hun badminton?} & \textit{Jeg vil gerne vide, \da{hvornår} \dabold{hun går} til badminton.} & \textit{Hun spiller hver torsdag aften fra kl. 19.00 til 20.30.} \\
\midrule
6. alene / sammen med & \textit{Spiller hun alene eller sammen med nogen?} & \textit{Kan du fortælle mig, \da{hvem} \dabold{hun spiller} sammen med?} & \textit{Hun spiller sammen med sin søster og to veninder.} \\
\end{longtable}
\normalsize

\begin{emailbox}{Opgave 2 第三人称完整信息汇总转述范文（82 词｜练习口语连贯总结）}
Sofie arbejder i Børnehaven Solstrålen i Odense. Hun arbejder på fuldtid fra mandag til fredag mellem klokken otte og halv fire, så hun har altid fri i weekenden. Hendes vigtigste arbejdsopgaver er at lege med børnene, læse bøger højt og gå tur med dem på legepladsen. I sin fritid kan hun rigtig godt lide at holde sig aktiv. Derfor spiller hun badminton i en lokal klub hver torsdag aften fra klokken syv til halv ni sammen med sin søster og to veninder.
\end{emailbox}
